\subsubsection{\texorpdfstring{$S_4$}{S4} 正则闭包的整体谱：属数、Chevalley--Weil 分解、正规子群商与 Jacobian 表示分裂}\label{subsubsec:xi-s4-regular-closure-global-spectrum}
沿用
\[
\Pi(\lambda,y)=\lambda^{4}-\lambda^{3}-(2y+1)\lambda^{2}+\lambda+y(y+1)\in\QQ(y)[\lambda],
\]
并记
\[
g(y):=256y^{3}+411y^{2}+165y+32,\qquad
\Delta(y):=-y(y-1)g(y).
\]
设 \(L/\QQ(y)\) 为 \(\Pi\) 的分裂域，\(G:=\mathrm{Gal}(L/\QQ(y))\cong S_4\)，并令 \(X\) 为 \(\QQ(X)=L\) 的光滑射影模型（即 \(X=\widetilde X\)）。

\begin{conclusion}[\texorpdfstring{$S_4$}{S4} 正则闭包曲线的精确属数]\label{con:xi-terminal-zm-s4-regular-closure-exact-genus-16}
覆盖 \(X\to\PP^1_y\) 为次数 \(24\) 的正则伽罗瓦覆盖，且
\[
g(X)=16.
\]
\end{conclusion}
\begin{proof}[证明要点]
由结论 \ref{con:xi-terminal-zm-s4-galois-closure-genus-tower-holomorphic-differentials}，分歧护照为一个阶 \(4\) 惯性点与五个阶 \(2\) 惯性点，即 \((4)\cdot(2)^5\)。
对正则伽罗瓦覆盖应用 Hurwitz 公式
\[
g(X)=1+\frac{|S_4|}{2}\!\left(-2+\sum_{p\in B}\Bigl(1-\frac1{e_p}\Bigr)\right),
\]
代入 \(|S_4|=24\)、\((e_p)=(4,2,2,2,2,2)\) 即得
\[
g(X)=1+12\!\left(-2+\frac34+5\cdot\frac12\right)=16.
\]
证毕。
\end{proof}

\begin{conclusion}[\texorpdfstring{$S_4$}{S4} 正则闭包中共轭类定点计数与循环商曲线属数]\label{con:xi-terminal-zm-s4-conjugacy-class-fixedpoints-and-cyclic-quotient-genera}
设
\(
\varphi:X\to \PP^1_y
\)
为结论 \ref{con:xi-terminal-zm-s4-regular-closure-exact-genus-16} 中的正则 \(S_4\)-覆盖，其分歧惯性型为一个 \(C_4\) 与五个 \(C_2\)。
对共轭类代表 \(g\in S_4\)，记 \(\mathrm{Fix}_X(g)\) 为其在 \(X\) 上的定点集，则有
\[
\begin{array}{c|c|c}
g\text{ 的循环型} & \#\mathrm{Fix}_X(g) & g(X/\langle g\rangle)\\
\hline
(2) & 10 & 6\\
(2,2) & 2 & 8\\
(4) & 2 & 4\\
(3) & 0 & 6
\end{array}
\]
\end{conclusion}
\begin{proof}[证明]
设 \(b\in \PP^1_y\) 为分歧值，\(I_b\le S_4\) 为其惯性群。由正则覆盖的陪集描述，
\(\varphi^{-1}(b)\cong S_4/I_b\)，且 \(xI_b\) 对应点的稳定子为 \(xI_bx^{-1}\)。故对任意 \(g\in S_4\) 有
\begin{equation}\label{eq:xi-terminal-zm-s4-fixedpoint-fiber-coset-formula}
\#\bigl(\mathrm{Fix}_X(g)\cap \varphi^{-1}(b)\bigr)
=\frac{\#\{x\in S_4\mid x^{-1}gx\in I_b\}}{|I_b|}.
\end{equation}

先取换位 \(\tau\)。五个有限分歧点 \(b_1,\dots,b_5\) 的惯性均为 \(C_2\)，其非平凡元均为换位。
由 \eqref{eq:xi-terminal-zm-s4-fixedpoint-fiber-coset-formula}，每个 \(b_i\) 的贡献为
\[
\frac{\#\{x\in S_4\mid x^{-1}\tau x=\tau_i\}}{2}
=\frac{|C_{S_4}(\tau)|}{2}
=\frac{4}{2}=2,
\]
故 \(\#\mathrm{Fix}_X(\tau)=10\)。对阶 \(2\) 商映射 \(X\to X/\langle\tau\rangle\) 用 Hurwitz 公式并代入 \(g(X)=16\)：
\[
2g(X)-2=2\bigl(2g(X/\langle\tau\rangle)-2\bigr)+\#\mathrm{Fix}_X(\tau),
\]
得 \(30=4g'-4+10\)，故 \(g'=6\)。

再取双换位 \(\delta\)。有限分歧惯性不含双换位，故仅无穷远分歧值 \(b_\infty\) 可能贡献。
记其惯性 \(I_{b_\infty}=\langle c\rangle\cong C_4\)，则唯一阶 \(2\) 元 \(c^2\) 为双换位。
由 \eqref{eq:xi-terminal-zm-s4-fixedpoint-fiber-coset-formula}
\[
\#\mathrm{Fix}_X(\delta)
=\frac{\#\{x\in S_4\mid x^{-1}\delta x=c^2\}}{4}
=\frac{|C_{S_4}(\delta)|}{4}
=\frac{8}{4}=2.
\]
再由 Hurwitz（阶 \(2\)）得
\[
30=2\bigl(2g(X/\langle\delta\rangle)-2\bigr)+2,
\]
故 \(g(X/\langle\delta\rangle)=8\)。

取四循环 \(\gamma\)。同理仅 \(b_\infty\) 贡献。
由于四循环同属一共轭类，可不失一般性取 \(\gamma=c\)。
此时式 \eqref{eq:xi-terminal-zm-s4-fixedpoint-fiber-coset-formula} 的条件 \(x^{-1}\gamma x\in \langle\gamma\rangle\) 等价于 \(x\in N_{S_4}(\langle\gamma\rangle)\)，
而 \(|N_{S_4}(\langle\gamma\rangle)|=8\)，故
\[
\#\mathrm{Fix}_X(\gamma)=\frac{8}{4}=2.
\]
又 \(\gamma^2\) 为双换位，已知 \(\#\mathrm{Fix}_X(\gamma^2)=2\)，且 \(\mathrm{Fix}_X(\gamma)\subseteq \mathrm{Fix}_X(\gamma^2)\)，
从而两者相等，不存在“仅被 \(\gamma^2\) 固定”之点。
故 \(X\to X/\langle\gamma\rangle\) 的分歧仅来自这 \(2\) 个点，且各分歧指数为 \(4\)。Hurwitz（阶 \(4\)）给出
\[
30=4\bigl(2g(X/\langle\gamma\rangle)-2\bigr)+2(4-1),
\]
解得 \(g(X/\langle\gamma\rangle)=4\)。

最后取三循环 \(\sigma\)。各惯性群仅含阶 \(2\) 或 \(4\) 元，故 \(\sigma\) 不可能固定分歧纤维点，从而 \(\mathrm{Fix}_X(\sigma)=\varnothing\)。
于是 \(X\to X/\langle\sigma\rangle\) 无分歧，Hurwitz（阶 \(3\)）给出
\[
30=3\bigl(2g(X/\langle\sigma\rangle)-2\bigr),
\]
故 \(g(X/\langle\sigma\rangle)=6\)。证毕。
\end{proof}

\begin{conclusion}[四循环商的 Prym 三维因子与 \texorpdfstring{$9$}{9} 维等型分量识别]\label{con:xi-terminal-zm-s4-c4-prym-threefold-and-v3sgn-identification}
取 \(c\in S_4\) 为无穷远惯性生成元（四循环），定义
\[
Y:=X/\langle c\rangle,\qquad
E_{\mathrm D}:=X/D_8
\]
（其中 \(D_8=N_{S_4}(\langle c\rangle)\)）。令
\[
\pi:Y\longrightarrow E_{\mathrm D}
\]
为诱导的次数 \(2\) 映射，并定义
\[
B:=\mathrm{Prym}(Y/E_{\mathrm D})
=\ker\!\bigl(\mathrm{Nm}_\pi:\mathrm{Jac}(Y)\to \mathrm{Jac}(E_{\mathrm D})\bigr)^0.
\]
则有：
\begin{enumerate}
  \item \(\pi\) 为分歧双覆盖，且 \(\deg\mathrm{Branch}(\pi)=6\)，从而 \(\dim B=3\)；
  \item \(\mathrm{Jac}(Y)\sim_{\QQ}E_{\mathrm D}\times B\)；
  \item 对任意换位 \(\tau\) 与双换位 \(\delta\)，有
  \[
  \mathrm{Jac}(X/\langle\tau\rangle)\sim_{\QQ}E^2\times E_{\mathrm D}\times B,
  \]
  \[
  \mathrm{Jac}(X/\langle\delta\rangle)\sim_{\QQ}\mathrm{Jac}(X_A)\times E\times E_{\mathrm D}^2\times B;
  \]
  \item 若 \(A_{V_3\otimes\mathrm{sgn}}\subset \mathrm{Jac}(X)\) 表示对应 \(V_3\otimes\mathrm{sgn}\) 的 \(9\) 维等型分量，则
  \[
  A_{V_3\otimes\mathrm{sgn}}\sim_{\QQ}B^3.
  \]
\end{enumerate}
\end{conclusion}
\begin{proof}[证明]
由子群塔 \(\langle c\rangle\subset D_8\subset S_4\) 可知 \(\deg(\pi)=2\)。
又由结论 \ref{con:xi-terminal-zm-s4-conjugacy-class-fixedpoints-and-cyclic-quotient-genera} 与
\ref{con:xi-terminal-zm-s4-hurwitz-subgroup-spectrum-cycletypes}，
\[
g(Y)=g(X/\langle c\rangle)=4,\qquad g(E_{\mathrm D})=g(X/D_8)=1.
\]
Hurwitz 公式
\[
2g(Y)-2=2\bigl(2g(E_{\mathrm D})-2\bigr)+\deg\mathrm{Branch}(\pi)
\]
给出 \(\deg\mathrm{Branch}(\pi)=6\)，并有
\(
\dim B=g(Y)-g(E_{\mathrm D})=3
\)。
这证明第 \(1\) 点。

第 \(2\) 点由双覆盖的 Prym 标准分解
\[
\mathrm{Jac}(Y)\sim_{\QQ}\pi^*\mathrm{Jac}(E_{\mathrm D})\times \mathrm{Prym}(Y/E_{\mathrm D})
\]
直接得到。

第 \(3\) 点与结论
\ref{con:xi-terminal-zm-s4-transposition-quotient-jacobian-full-decomposition}、
\ref{con:xi-terminal-zm-s4-doubletransposition-quotient-jacobian-full-decomposition}
一致：其中文中记号 \(E_{\mathrm{res}}\) 即此处 \(E_{\mathrm D}\)，且
\(
J_{\mathrm{LY}}\sim_{\QQ}\mathrm{Jac}(X_A)
\)
见结论
\ref{con:xi-terminal-zm-s4-quotient-genus-spectrum-and-jacobian-representation-splitting}。

第 \(4\) 点可由两种等价方式得到：其一，由结论
\ref{con:xi-terminal-zm-s4-jacobian-group-algebra-prym-e3-b3}
的分解
\[
\mathrm{Jac}(X)\sim_{\QQ}\mathrm{Jac}(X_A)\times E_{\mathrm D}^2\times E^3\times B^3
\]
直接识别 \(V_3\otimes\mathrm{sgn}\) 等型即为 \(B^3\)；
其二，由
\(
e_{V_3\otimes\mathrm{sgn}}\QQ[S_4]e_{V_3\otimes\mathrm{sgn}}\cong M_3(\QQ)
\)
与原始幂等分解得到同结论。证毕。
\end{proof}

\begin{conclusion}[Hasse--Weil 动机分解与良素处局部 \texorpdfstring{$L$}{L} 因子幂次刚性]\label{con:xi-terminal-zm-s4-hasse-weil-full-factorization-rigid-exponents}
在 \(\QQ\)-同源意义下有
\[
\mathrm{Jac}(X)\sim_{\QQ}\mathrm{Jac}(X_A)\times E^3\times E_{\mathrm D}^2\times B^3.
\]
因而对任一使上述各因子皆良约化的素数 \(p\)，局部 \(L\)-多项式分子满足
\[
P_p(X;T)
=P_p(X_A;T)\cdot P_p(E;T)^3\cdot P_p(E_{\mathrm D};T)^2\cdot P_p(B;T)^3,
\]
并且
\(
\deg P_p(B;T)=2\dim B=6
\)。
\end{conclusion}
\begin{proof}[证明]
第一式由结论
\ref{con:xi-terminal-zm-s4-c4-prym-threefold-and-v3sgn-identification}
第 \(4\) 点与既有 \(A_{\mathrm{sgn}}\)、\(A_{V_3}\)、\(A_{V_2}\) 识别合并得到。
第二式由“同源阿贝尔簇在良素处具有相同局部 \(L\)-因子”及乘积阿贝尔簇局部因子的乘法性直接推出。证毕。
\end{proof}

\begin{conclusion}[全纯 \(1\)-形式的 Chevalley--Weil 刚性分解]\label{con:xi-terminal-zm-s4-holomorphic-oneforms-chevalley-weil-rigid-decomposition}
记 \(S_4\) 的不可约表示为
\[
\mathbf 1,\ \mathrm{sgn},\ V_3,\ V_3\otimes\mathrm{sgn},\ V_2.
\]
则有
\[
H^0(X,\Omega_X^1)\cong
\mathrm{sgn}^{\oplus 2}\oplus V_3\oplus (V_3\otimes\mathrm{sgn})^{\oplus 3}\oplus V_2,
\]
并且平凡表示不出现。
\end{conclusion}
\begin{proof}[证明要点]
对 \(\rho\neq \mathbf 1\)，Chevalley--Weil 公式在本情形可写为
\[
m_\rho=-\dim\rho+\frac12\sum_{p\in B}\bigl(\dim\rho-\dim(\rho^{I_p})\bigr),
\]
且 \(m_{\mathbf 1}=0\)（底曲线 \(\PP^1\) 无全纯 \(1\)-形式）。
本覆盖仅需 \(C_2\) 与 \(C_4\) 两类惯性不变维数：
\[
\begin{aligned}
&\dim(\mathrm{sgn}^{C_2})=0,\ \dim(\mathrm{sgn}^{C_4})=0,\\
&\dim(V_3^{C_2})=2,\ \dim(V_3^{C_4})=0,\\
&\dim((V_3\!\otimes\!\mathrm{sgn})^{C_2})=1,\ \dim((V_3\!\otimes\!\mathrm{sgn})^{C_4})=1,\\
&\dim(V_2^{C_2})=1,\ \dim(V_2^{C_4})=1.
\end{aligned}
\]
代入五个 \(C_2\) 分歧点与一个 \(C_4\) 分歧点，得到
\[
m_{\mathrm{sgn}}=2,\quad m_{V_3}=1,\quad m_{V_3\otimes\mathrm{sgn}}=3,\quad m_{V_2}=1,\quad m_{\mathbf 1}=0.
\]
与结论 \ref{con:xi-terminal-zm-s4-chevalley-weil-explicit-decomposition} 一致。证毕。
\end{proof}

\begin{conclusion}[正规子群商曲线的显式模型与属数：\texorpdfstring{$A_4$}{A4} 商与 \texorpdfstring{$V_4$}{V4} 商]\label{con:xi-terminal-zm-s4-normal-subgroup-quotients-explicit-models-genus}
设
\[
X_A:=X/A_4,\qquad X_V:=X/V_4^{\mathrm{norm}}.
\]
则：
\begin{enumerate}
  \item \(X_A\to\PP^1_y\) 为二次覆盖，且
  \[
  \QQ(X_A)=\QQ\!\left(y,\sqrt{\Delta(y)}\right),
  \qquad
  X_A:\ t^2=\Delta(y),
  \qquad g(X_A)=2.
  \]
  \item \(X_V\to\PP^1_y\) 为次数 \(6\) 的 \(S_3\)-正则覆盖，且
  \[
  \QQ(X_V)=\QQ(y,z,\sqrt{\Delta(y)}),\qquad \mathscr R(z,y)=0,
  \]
  其中可取
  \[
  \mathscr R(z,y)=z^3+(2y+1)z^2-(4y^2+4y+1)z-(8y^3+13y^2+5y+1),
  \]
  并且 \(g(X_V)=4\)。
\end{enumerate}
\end{conclusion}
\begin{proof}[证明要点]
对 \(A_4\)：指数 \(2\) 正规子群对应符号同态核，故固定域由判别式平方根生成，即
\(
\QQ(X_A)=\QQ(y,\sqrt{\Delta(y)})
\)。
由于 \(\deg\Delta=5\) 为奇数，双覆盖在五个有限零点与无穷远点分歧，总分歧点数 \(6\)，故 \(g(X_A)=2\)。
\par
对 \(V_4^{\mathrm{norm}}\)：商群 \(S_4/V_4^{\mathrm{norm}}\cong S_3\)，固定域由三次解式闭包给出。
以 \(\mathscr R\) 为 resolvent cubic，附加其判别式平方根即得 \(S_3\)-闭包域，
从而 \(\QQ(X_V)=\QQ(y,z,\sqrt{\Delta(y)})\)。
由结论 \ref{con:xi-terminal-zm-s4-hurwitz-subgroup-spectrum-cycletypes} 已知 \(g(X_V)=4\)；等价地，用 \(|S_3|=6\) 与六个阶 \(2\) 分歧值代入 Hurwitz 公式亦得同值。证毕。
\end{proof}

\begin{conclusion}[双换位商曲线的双二次结构与纤维积刻画]\label{con:xi-terminal-zm-s4-delta-quotient-v4-fiberproduct}
沿用本小节记号，取无穷远惯性生成元 \(c\in S_4\) 为四循环，置
\[
\delta:=c^2,\qquad
Z:=X/\langle\delta\rangle,\qquad
E_{\mathrm D}:=X/D_8,\qquad
Y:=X/\langle c\rangle.
\]
则有：
\begin{enumerate}
  \item \(g(Z)=8\)；
  \item \(Z\to E_{\mathrm D}\) 为次数 \(4\) 的 Galois 覆盖，覆盖群同构于 \(V_4\)；
  \item 在函数域层面
  \[
  \QQ(Z)=\QQ(X_V)\cdot\QQ(Y)\quad\text{于}\ \QQ(E_{\mathrm D})\ \text{之上},
  \]
  且
  \[
  \QQ(X_V)\cap\QQ(Y)=\QQ(E_{\mathrm D}),
  \]
  从而 \(Z\) 同构于纤维积 \(X_V\times_{E_{\mathrm D}}Y\) 的规范化。
\end{enumerate}
\end{conclusion}
\begin{proof}
第 \(1\) 点由结论 \ref{con:xi-terminal-zm-s4-conjugacy-class-fixedpoints-and-cyclic-quotient-genera} 中
\(
\#\mathrm{Fix}_X(\delta)=2
\)
与 \(g(X)=16\) 代入二次商 Hurwitz 公式：
\[
2g(X)-2=2(2g(Z)-2)+\#\mathrm{Fix}_X(\delta)
\]
即得 \(g(Z)=8\)。

第 \(2\) 点中，\(\delta\in\langle c\rangle\subset D_8\)，且 \(V_4^{\mathrm{norm}}\subset D_8\)。
并有
\[
\langle\delta\rangle=V_4^{\mathrm{norm}}\cap\langle c\rangle,\qquad
D_8/\langle\delta\rangle\cong V_4.
\]
故 \(Z=X/\langle\delta\rangle\to X/D_8=E_{\mathrm D}\) 为次数 \(4\) 的 \(V_4\)-Galois 覆盖。

第 \(3\) 点令 \(L:=\QQ(X)\)，则
\[
\QQ(X_V)=L^{V_4^{\mathrm{norm}}},\qquad
\QQ(Y)=L^{\langle c\rangle},\qquad
\QQ(E_{\mathrm D})=L^{D_8}.
\]
由 Galois 对应性，
\[
L^{V_4^{\mathrm{norm}}}\cap L^{\langle c\rangle}
=L^{\langle V_4^{\mathrm{norm}},\langle c\rangle\rangle}
=L^{D_8}
=\QQ(E_{\mathrm D}),
\]
\[
L^{V_4^{\mathrm{norm}}}\cdot L^{\langle c\rangle}
=L^{V_4^{\mathrm{norm}}\cap\langle c\rangle}
=L^{\langle\delta\rangle}
=\QQ(Z).
\]
由“纤维积规范化的函数域等于子域合成”即得结论。证毕。
\end{proof}

\begin{conclusion}[三条二次子覆盖的分支奇偶对称差律]\label{cor:xi-terminal-zm-s4-v4-three-quad-subcovers-branch-parity}
在结论 \ref{con:xi-terminal-zm-s4-delta-quotient-v4-fiberproduct} 的 \(V_4\)-Galois 覆盖
\(
Z\to E_{\mathrm D}
\)
中，存在且仅存在三条非平凡二次子覆盖：
\[
X_V\to E_{\mathrm D},\qquad
Y\to E_{\mathrm D},\qquad
X_H\to E_{\mathrm D},
\]
其中可取 \(H=\langle\delta,\tau\rangle\le D_8\)（\(\tau\) 为 \(D_8\) 内换位），并置 \(X_H:=X/H\)。则：
\begin{enumerate}
  \item \(g(X_H)=2\)，且 \(X_H\to E_{\mathrm D}\) 为分支次数 \(2\) 的二次覆盖；
  \item 若
  \[
  \QQ(X_V)=\QQ(E_{\mathrm D})(\sqrt{f_V}),\qquad
  \QQ(Y)=\QQ(E_{\mathrm D})(\sqrt{f_Y}),
  \]
  则
  \[
  \QQ(X_H)=\QQ(E_{\mathrm D})(\sqrt{f_Vf_Y}),
  \]
  并满足
  \[
  \mathrm{Branch}(X_H/E_{\mathrm D})
  \equiv
  \mathrm{Branch}(X_V/E_{\mathrm D})\triangle \mathrm{Branch}(Y/E_{\mathrm D})
  \pmod 2.
  \]
  \item 设 \(D_V,D_Y,D_H\) 为三条二次覆盖的约化分支集合。
  已知 \(\deg D_V=\deg D_Y=6,\ \deg D_H=2\)，故
  \[
  \deg(D_V\cap D_Y)=5,
  \]
  即 \(D_V,D_Y\) 共有 \(5\) 个分支点，且各自仅有 \(1\) 个独有分支点；这两个独有点恰构成 \(D_H\)。
\end{enumerate}
\end{conclusion}
\begin{proof}
对第 \(1\) 点，取 \(H=\langle\delta,\tau\rangle=\{1,\delta,\tau,\delta\tau\}\)。
由结论 \ref{con:xi-terminal-zm-s4-conjugacy-class-fixedpoints-and-cyclic-quotient-genera}，
\[
\#\mathrm{Fix}_X(\delta)=2,\qquad
\#\mathrm{Fix}_X(\tau)=\#\mathrm{Fix}_X(\delta\tau)=10.
\]
群作用 Hurwitz 公式给出
\[
2g(X)-2=|H|\bigl(2g(X_H)-2\bigr)+\sum_{1\neq h\in H}\#\mathrm{Fix}_X(h),
\]
即
\[
30=4\bigl(2g(X_H)-2\bigr)+(2+10+10),
\]
故 \(g(X_H)=2\)。
又 \([D_8:H]=2\)，故 \(X_H\to E_{\mathrm D}\) 为二次覆盖；由 \(g(E_{\mathrm D})=1\) 得分支次数 \(2g(X_H)-2=2\)。

第 \(2\) 点是双二次扩张的局部赋值判别：在任意离散赋值 \(v\) 下，\(\QQ(E_{\mathrm D})(\sqrt{f})/\QQ(E_{\mathrm D})\) 分歧当且仅当 \(v(f)\) 为奇数。对 \(f_V,f_Y,f_Vf_Y\) 逐点比较即得模 \(2\) 对称差律。

第 \(3\) 点由第 \(1\) 点与已有 \(g(X_V)=g(Y)=4\) 立即得到 \(\deg D_V=\deg D_Y=6,\deg D_H=2\)。
在特征 \(0\) 下二次分支为简单分支，故
\[
2=\deg(D_V\triangle D_Y)=\deg D_V+\deg D_Y-2\deg(D_V\cap D_Y)=12-2\deg(D_V\cap D_Y),
\]
从而 \(\deg(D_V\cap D_Y)=5\)。证毕。
\end{proof}

\begin{conclusion}[\(\delta\)-不变/反不变分解与 \(\mathrm{Jac}(Z)\)、\(\mathrm{Prym}(X/Z)\) 同源型]\label{con:xi-terminal-zm-s4-delta-eigensplitting-jacz-prymxz}
在结论 \ref{con:xi-terminal-zm-s4-delta-quotient-v4-fiberproduct} 的记号下，存在定义于 \(\QQ\) 上的同源分解
\[
\mathrm{Jac}(Z)\sim_{\QQ}\mathrm{Jac}(X_A)\times E\times E_{\mathrm D}^2\times B,
\]
以及
\[
\mathrm{Prym}(X/Z)\sim_{\QQ}E^2\times B^2.
\]
\end{conclusion}
\begin{proof}
记 \(p:X\to Z=X/\langle\delta\rangle\)，并在 \(\QQ[\langle\delta\rangle]\) 中取幂等元
\[
e_+:=\frac{1+\delta}{2},\qquad e_-:=\frac{1-\delta}{2}.
\]
标准理论给出
\[
\mathrm{Jac}(X)\sim_{\QQ}e_+\mathrm{Jac}(X)\times e_-\mathrm{Jac}(X),
\]
且
\(
e_+\mathrm{Jac}(X)\sim_{\QQ}\mathrm{Jac}(Z)
\),
\(
e_-\mathrm{Jac}(X)\sim_{\QQ}\mathrm{Prym}(X/Z)
\)。

再由结论 \ref{con:xi-terminal-zm-s4-jacobian-group-algebra-prym-e3-b3}
\[
\mathrm{Jac}(X)\sim_{\QQ}\mathrm{Jac}(X_A)\times E^3\times E_{\mathrm D}^2\times B^3.
\]
对双换位 \(\delta\)，其在 \(\mathrm{sgn},V_2\) 上取值为 \(+1\)，在 \(V_3,V_3\otimes\mathrm{sgn}\) 上迹为 \(-1\)，故 \(+1\)-特征维数分别为
\[
1,\ 2,\ 1,\ 1.
\]
于是 \(e_+\)-像分别贡献
\[
\mathrm{Jac}(X_A),\ E_{\mathrm D}^2,\ E,\ B,
\]
而 \(e_-\)-像分别贡献
\[
0,\ 0,\ E^2,\ B^2.
\]
合并即得结论。证毕。
\end{proof}

\begin{conclusion}[属 \(2\) 中间商曲线的 Jacobian 分裂]\label{con:xi-terminal-zm-s4-xh-jacobian-splitting-ed-e}
取
\[
H=\langle\delta,\tau\rangle\le D_8,\qquad X_H:=X/H
\]
（\(\tau\) 为 \(D_8\) 中换位）。则 \(g(X_H)=2\)，且
\[
\mathrm{Jac}(X_H)\sim_{\QQ}E\times E_{\mathrm D}.
\]
特别地，\(X_H\to E_{\mathrm D}\) 为分支次数 \(2\) 的二次覆盖，且 \(X_H\) 为双椭圆型属 \(2\) 曲线。
\end{conclusion}
\begin{proof}
由结论 \ref{cor:xi-terminal-zm-s4-v4-three-quad-subcovers-branch-parity} 已知 \(g(X_H)=2\)。
令
\[
e_H:=\frac{1}{|H|}\sum_{h\in H}h\in\QQ[H].
\]
则 \(e_H\mathrm{Jac}(X)\sim_{\QQ}\mathrm{Jac}(X_H)\)。
对任一不可约表示 \(\rho\)，有
\[
\dim(\rho^H)=\frac14\!\left(\chi_\rho(1)+\chi_\rho(\delta)+2\chi_\rho(\text{换位})\right).
\]
代入 \(S_4\) 特征标值得
\[
\dim(\mathrm{sgn}^H)=0,\quad
\dim((V_3\otimes\mathrm{sgn})^H)=0,\quad
\dim(V_3^H)=1,\quad
\dim(V_2^H)=1.
\]
故 \(\mathrm{Jac}(X_A)\) 与 \(B^3\) 不贡献不变部分，而 \(E^3\) 与 \(E_{\mathrm D}^2\) 各贡献一维椭圆因子，遂有
\[
\mathrm{Jac}(X_H)\sim_{\QQ}E\times E_{\mathrm D}.
\]
其余断言由 \(g(X_H)=2\) 与该分裂立即得到。证毕。
\end{proof}

\begin{conclusion}[三次解式判别式与 \texorpdfstring{$\Delta(y)$}{Delta(y)} 的一致性]\label{con:xi-terminal-zm-s4-resolvent-cubic-discriminant-equals-delta}
在结论 \ref{con:xi-terminal-zm-s4-normal-subgroup-quotients-explicit-models-genus} 的记号下，恒有
\[
\mathrm{Disc}_z\!\bigl(\mathscr R(z,y)\bigr)=\Delta(y)=-y(y-1)g(y).
\]
\end{conclusion}
\begin{proof}[证明要点]
将 \(\mathscr R=z^3+az^2+bz+c\) 代入三次判别式闭式
\[
\mathrm{Disc}_z(\mathscr R)=a^2b^2-4b^3-4a^3c-27c^2+18abc,
\]
其中
\[
a=2y+1,\quad b=-(4y^2+4y+1),\quad c=-(8y^3+13y^2+5y+1),
\]
直接化简即得 \(-y(y-1)g(y)\)。
该恒等式亦与结论 \ref{con:xi-terminal-zm-discriminant-leyang}、\ref{con:xi-terminal-zm-s4-c3-fiberproduct-normalization} 的二次分辨数据一致。证毕。
\end{proof}

\begin{conclusion}[中间商曲线属数谱与 Jacobian 的表示论分裂]\label{con:xi-terminal-zm-s4-quotient-genus-spectrum-and-jacobian-representation-splitting}
对任意子群 \(H\le S_4\)，有
\[
g(X/H)=\dim_{\CC}H^0(X,\Omega_X^1)^H.
\]
由结论 \ref{con:xi-terminal-zm-s4-holomorphic-oneforms-chevalley-weil-rigid-decomposition} 可直接得到
\[
g(X/A_4)=2,\quad
g(X/V_4^{\mathrm{norm}})=4,\quad
g(X/S_3)=1,\quad
g(X/C_3)=6,\quad
g(X/D_8)=1.
\]
此外，存在定义在 \(\QQ\) 上的阿贝尔簇
\[
A_{\mathrm{sgn}},\ A_{V_3},\ A_{V_3\otimes\mathrm{sgn}},\ A_{V_2},
\]
使
\[
\mathrm{Jac}(X)\sim_{\QQ}
A_{\mathrm{sgn}}\times A_{V_3}\times A_{V_3\otimes\mathrm{sgn}}\times A_{V_2},
\]
并且
\[
\dim A_{\mathrm{sgn}}=2,\quad
\dim A_{V_3}=3,\quad
\dim A_{V_3\otimes\mathrm{sgn}}=9,\quad
\dim A_{V_2}=2.
\]
同时有规范同源识别
\[
A_{\mathrm{sgn}}\sim_{\QQ}\mathrm{Jac}(X/A_4),\qquad
A_{\mathrm{sgn}}\times A_{V_2}\sim_{\QQ}\mathrm{Jac}(X/V_4^{\mathrm{norm}}).
\]
\end{conclusion}
\begin{proof}[证明要点]
恒等式 \(H^0(X/H,\Omega^1)\cong H^0(X,\Omega^1)^H\) 给出第一式。
将
\[
H^0(X,\Omega_X^1)\cong
2\,\mathrm{sgn}\oplus V_3\oplus 3(V_3\otimes\mathrm{sgn})\oplus V_2
\]
限制到各子群并取不变维数，即得所列属数（与结论
\ref{con:xi-terminal-zm-s4-hurwitz-subgroup-spectrum-cycletypes} 一致）。
\par
对 Jacobian，按 \(\QQ[S_4]\) 中心幂等元作用得到 isotypic 分裂；由于 \(S_4\) 不可约特征均为有理值，各因子可在 \(\QQ\) 上定义。
维数由重数乘表示维数给出：
\[
2\cdot 1,\ 1\cdot 3,\ 3\cdot 3,\ 1\cdot 2,
\]
即 \((2,3,9,2)\)。
当 \(H\) 正规（如 \(A_4,V_4^{\mathrm{norm}}\)）时，\(\mathrm{Jac}(X/H)\) 正对应于在 \(H\) 上平凡的表示分量之直和，故得两条规范同源识别。证毕。
\end{proof}
\paragraph{\texorpdfstring{$S_4$}{S4}-等变 \texorpdfstring{$H^1$}{H1} 分解、驯导子谱与函数域 Langlands 分量化}
沿用子小节 \ref{subsubsec:xi-s4-regular-closure-global-spectrum} 的记号，令
\(
X=\widetilde X
\)
为 Lee--Yang 四次谱覆盖的 \(S_4\)-正则闭包曲线，分歧惯性型为 \((4,2,2,2,2,2)\)。
记 \(S_4\) 的四个非平凡不可约表示为
\[
\varepsilon=\mathrm{sgn},\qquad V_2,\qquad V_3,\qquad V_3':=V_3\otimes\varepsilon.
\]

\begin{conclusion}[等变 Weil--Artin 分量分解的次数与幂指数刚性]\label{con:xi-terminal-zm-s4-equivariant-weil-artin-isotypic-factorization}
设 \(p\nmid 2\cdot 3\cdot 31\cdot 37\)，并设 \(X\) 在 \(p\) 处具有光滑良约化 \(X_p/\FF_p\)。
取素数 \(\ell\neq p\)，并记几何 Frobenius 为 \(\mathrm{Frob}_p\)。
则 \(H^1_{\acute et}(X_{\overline{\FF}_p},\QQ_\ell)\) 作为 \(\QQ_\ell[S_4]\)-半单模存在唯一的等变直和分解
\[
H^1_{\acute et}(X_{\overline{\FF}_p},\QQ_\ell)
\ \cong\
(\varepsilon\otimes M_{\varepsilon})
\ \oplus\
(V_2\otimes M_{2})
\ \oplus\
(V_3\otimes M_{3})
\ \oplus\
(V_3'\otimes M_{3'}),
\]
其中
\(
M_\sigma:=\Hom_{S_4}\!\bigl(\sigma,H^1_{\acute et}(X_{\overline{\FF}_p},\QQ_\ell)\bigr)
\)
为 \(\mathrm{Gal}(\overline{\FF}_p/\FF_p)\)-表示，并满足
\[
\dim M_{\varepsilon}=4,\qquad
\dim M_{2}=2,\qquad
\dim M_{3}=2,\qquad
\dim M_{3'}=6.
\]
令
\[
P_\sigma(T):=\det\!\left(1-T\mathrm{Frob}_p\mid M_\sigma\right)\in 1+T\overline{\QQ}[T]
\qquad(\sigma\in\{\varepsilon,2,3,3'\}),
\]
并记曲线的局部因子分子
\(
P_p(X;T):=\det(1-T\mathrm{Frob}_p\mid H^1_{\acute et}(X_{\overline{\FF}_p},\QQ_\ell))
\)。
则有幂指数分解
\[
P_p(X;T)=P_{\varepsilon}(T)^{1}\cdot P_{2}(T)^{2}\cdot P_{3}(T)^{3}\cdot P_{3'}(T)^{3},
\]
并且
\[
\deg P_{\varepsilon}=4,\qquad \deg P_{2}=2,\qquad \deg P_{3}=2,\qquad \deg P_{3'}=6.
\]
因此
\[
Z(X_p,T)=\frac{P_{\varepsilon}(T)\,P_{2}(T)^{2}\,P_{3}(T)^{3}\,P_{3'}(T)^{3}}{(1-T)(1-pT)}.
\]
\end{conclusion}
\begin{proof}
对任意 \(\ell\neq p\)，有自然同构
\(
H^1_{\acute et}(X_{\overline{\FF}_p},\QQ_\ell)\cong V_\ell(\mathrm{Jac}(X))
\)
且该同构与 \(S_4\) 作用及 \(\mathrm{Frob}_p\) 作用相容。
由结论 \ref{con:xi-terminal-zm-s4-quotient-genus-spectrum-and-jacobian-representation-splitting}，
\(\mathrm{Jac}(X)\) 在 \(\QQ\)-同源意义下按 \(\varepsilon,V_2,V_3,V_3'\) 的等型分量分裂，并满足
\[
\dim A_{\varepsilon}=2,\qquad
\dim A_{V_2}=2,\qquad
\dim A_{V_3}=3,\qquad
\dim A_{V_3'}=9.
\]
因此 \(H^1_{\acute et}\) 的相应等型分量维数分别为 \(4,4,6,18\)，再除以 \(\dim\sigma\in\{1,2,3,3\}\) 得
\[
\dim M_{\varepsilon}=4,\qquad
\dim M_{2}=2,\qquad
\dim M_{3}=2,\qquad
\dim M_{3'}=6,
\]
从而得到所示等变分解与次数数据。
再由 Schur 引理，
\(
\det(1-T\mathrm{Frob}_p\mid \sigma\otimes M_\sigma)=\det(1-T\mathrm{Frob}_p\mid M_\sigma)^{\dim\sigma}
\)，
故 \(P_p(X;T)=\prod_\sigma P_\sigma(T)^{\dim\sigma}\)。
次数陈述由 \(\deg P_\sigma=\dim M_\sigma\) 得出；最后一式为曲线的标准 ζ-函数表达式。证毕。
\end{proof}

\begin{conclusion}[分歧集上的驯 Artin 导子谱与全局导子度]\label{con:xi-terminal-zm-s4-tame-artin-conductor-spectrum}
设
\[
U:=\PP^1_y\setminus\Bigl(\{0,1,\infty\}\cup\{P_{\mathrm{LY}}(y)=0\}\Bigr),
\]
并令 \(\pi_1(U)\twoheadrightarrow S_4\) 为由正则闭包 \(X\to\PP^1_y\) 给出的单值化满射。
对每个不可约表示 \(\sigma\in\{\varepsilon,V_2,V_3,V_3'\}\)，记相应驯 lisse 局部系统为 \(\mathcal L_\sigma\)。
则对每个分歧点 \(b\)（其惰性群 \(I_b\) 为 \(C_2\) 或 \(C_4\)）有驯 Artin 导子指数
\[
a_b(\sigma)=\dim\sigma-\dim\sigma^{I_b},
\]
并且在五个有限分歧点（惰性生成元可取换位 \(\tau\)）处有
\[
a_b(\varepsilon)=1,\qquad a_b(V_2)=1,\qquad a_b(V_3)=1,\qquad a_b(V_3')=2,
\]
在无穷远分歧点（惰性生成元可取四循环 \(\rho\)）处有
\[
a_\infty(\varepsilon)=1,\qquad a_\infty(V_2)=1,\qquad a_\infty(V_3)=3,\qquad a_\infty(V_3')=2.
\]
因而四个表示的全局驯导子度
\(
A(\sigma):=\sum_b a_b(\sigma)
\)
分别为
\[
A(\varepsilon)=6,\qquad A(V_2)=6,\qquad A(V_3)=8,\qquad A(V_3')=12.
\]
\end{conclusion}
\begin{proof}
有限分歧点处 \(I_b\cong\langle\tau\rangle\)（\(\tau\) 为换位），无穷远处 \(I_\infty\cong\langle\rho\rangle\)（\(\rho\) 为四循环）。
对任意循环群 \(\langle g\rangle\)（阶 \(n\)）有角色平均
\[
\dim\sigma^{\langle g\rangle}=\frac1n\sum_{k=0}^{n-1}\chi_\sigma(g^k).
\]
取 \(S_4\) 角色值（按共轭类 \(1,(12),(12)(34),(123),(1234)\)）：
\[
\chi_{\varepsilon}=(1,-1,1,1,-1),\quad
\chi_{V_2}=(2,0,2,-1,0),\quad
\chi_{V_3}=(3,1,-1,0,-1),\quad
\chi_{V_3'}=(3,-1,-1,0,1).
\]
于是
\(
\dim\sigma^{\langle\tau\rangle}=\tfrac12(\chi_\sigma(1)+\chi_\sigma(\tau))
\)
给出
\[
\dim\varepsilon^{\langle\tau\rangle}=0,\quad
\dim V_2^{\langle\tau\rangle}=1,\quad
\dim V_3^{\langle\tau\rangle}=2,\quad
\dim V_3'^{\langle\tau\rangle}=1,
\]
从而五个有限分歧点的导子指数如表。
另一方面 \(\rho\) 为四循环且 \(\rho^2\) 为双换位，故
\(
\dim\sigma^{\langle\rho\rangle}=\tfrac14(\chi_\sigma(1)+2\chi_\sigma(\rho)+\chi_\sigma(\rho^2))
\)
给出
\[
\dim\varepsilon^{\langle\rho\rangle}=0,\quad
\dim V_2^{\langle\rho\rangle}=1,\quad
\dim V_3^{\langle\rho\rangle}=0,\quad
\dim V_3'^{\langle\rho\rangle}=1,
\]
从而无穷远处导子指数如表。
对六个分歧点求和即得全局导子度。证毕。
\end{proof}

\begin{conclusion}[函数域 Langlands 的分量化实现与导子精确性]\label{con:xi-terminal-zm-s4-function-field-langlands-artin-components}
设 \(p\nmid 2\cdot 3\cdot 31\cdot 37\)，并令 \(F=\FF_p(y)\)。
对每个不可约表示 \(\sigma\in\{\varepsilon,V_2,V_3,V_3'\}\)，考虑由正则闭包单值化得到的纯权 \(0\) 驯 lisse 局部系统 \(\mathcal L_\sigma\)（定义于 \(U\subset\PP^1_{\FF_p}\)）。
则存在一个（在同构意义下）唯一的尖点自守表示
\[
\pi_{\sigma,p}\ \subset\ \mathrm{GL}_{\dim\sigma}(\mathbb A_F),
\]
使得对每个非分歧处 \(v\) 的未分歧局部 \(L\)-因子有
\[
L_v(\pi_{\sigma,p},T)
=
\det\!\left(1-T\mathrm{Frob}_v\mid (\mathcal L_{\sigma,\bar\eta})^{I_v}\right)^{-1}.
\]
并且 \(\pi_{\sigma,p}\) 的算术导子在每个分歧点处的局部指数等于结论 \ref{con:xi-terminal-zm-s4-tame-artin-conductor-spectrum} 所给的 \(a_b(\sigma)\)，从而
\[
\deg\mathfrak f(\pi_{\varepsilon,p})=6,\qquad
\deg\mathfrak f(\pi_{V_2,p})=6,\qquad
\deg\mathfrak f(\pi_{V_3,p})=8,\qquad
\deg\mathfrak f(\pi_{V_3',p})=12.
\]
\end{conclusion}
\begin{proof}[证明要点]
由于 \(\mathcal L_\sigma\) 具有有限幺半单单值群，且仅在分歧集处驯分歧，函数域全球 Langlands 对应（Lafforgue）将其与唯一尖点自守表示 \(\pi_{\sigma,p}\) 配对，并逐点匹配未分歧局部 \(L\)-因子并保持导子指数。
将结论 \ref{con:xi-terminal-zm-s4-tame-artin-conductor-spectrum} 的驯导子谱代入即可得到导子度公式。证毕。\cite{LafforgueGlobalLanglandsFunctionFields}
\end{proof}

\begin{conclusion}[分量化 Weil 圆周性质与正号函数方程]\label{con:xi-terminal-zm-s4-componentwise-rh-functional-equation-plus}
在结论 \ref{con:xi-terminal-zm-s4-equivariant-weil-artin-isotypic-factorization} 的记号下，对每个 \(\sigma\in\{\varepsilon,2,3,3'\}\)：
\begin{enumerate}
  \item \(P_\sigma(T)\) 的全部逆根 \(\alpha\) 满足 \(|\alpha|=\sqrt p\)（任一复嵌入下），等价地其全部零点位于圆周 \(|T|=p^{-1/2}\)。
  \item 记 \(m_\varepsilon=2,m_2=1,m_3=1,m_{3'}=3\) 为 \(H^0(X,\Omega^1)\) 中的重数（结论 \ref{con:xi-terminal-zm-s4-holomorphic-oneforms-chevalley-weil-rigid-decomposition}）。
  则 \(\deg P_\sigma=2m_\sigma\)，并且函数方程具有正号形式
  \[
  P_\sigma(T)=p^{m_\sigma}\,T^{2m_\sigma}\,P_\sigma\!\left(\frac{1}{pT}\right).
  \]
\end{enumerate}
\end{conclusion}
\begin{proof}[证明要点]
纯性给出圆周性质：\(H^1_{\acute et}(X_{\overline{\FF}_p},\QQ_\ell)\) 为权 \(1\) 的纯表示，故其 Frobenius 特征值绝对值为 \(\sqrt p\)，并对任一 Frobenius 稳定子表示（尤其各 \(S_4\)-等型分量 \(\sigma\otimes M_\sigma\)）保持。
\par
函数方程来自 Poincar\'e 对偶：\(H^1\) 上存在 \(\mathrm{Frob}_p\)-半线性扭转的非退化交替配对，且 \(S_4\) 作用保持该配对。
由于四个 \(\sigma\) 皆自对偶且可取 \(S_4\)-不变对称双线性形式，配对在 \(\sigma\otimes M_\sigma\) 上诱导出 \(M_\sigma\) 的交替配对；于是 \(\det(1-T\mathrm{Frob}_p\mid M_\sigma)\) 为回文多项式且符号为 \(+1\)，指数由 \(\deg P_\sigma=\dim M_\sigma=2m_\sigma\) 给出。证毕。\cite{DeligneWeilII}
\end{proof}

\endinput


\paragraph{\texorpdfstring{$V_4$}{V4} 中间塔的分支刚化、幂等分解与 Prym 位置唯一化}

\begin{conclusion}[双二次覆盖 \texorpdfstring{$Z\to E_{\mathrm D}$}{Z to ED} 的分支总数与惯性分布]\label{con:xi-terminal-zm-s4-v4-ed-branchcount-inertia-115}
在结论 \ref{con:xi-terminal-zm-s4-delta-quotient-v4-fiberproduct} 与
\ref{cor:xi-terminal-zm-s4-v4-three-quad-subcovers-branch-parity} 的记号下，记
\[
\pi:Z\longrightarrow E_{\mathrm D},\qquad \operatorname{Gal}(Z/E_{\mathrm D})\cong V_4.
\]
则成立：
\begin{enumerate}
  \item \(\pi\) 在 \(E_{\mathrm D}\) 上恰有 \(7\) 个分支点，且均为惯性阶 \(2\) 的简单分支；
  \item 存在 \(V_4\setminus\{1\}=\{a,b,\gamma\}\) 的标号，使分支集合分解为互不相交约化有效除子
  \[
  D_a\sqcup D_b\sqcup D_\gamma,\qquad
  \deg D_a=\deg D_b=1,\ \deg D_\gamma=5,
  \]
  并满足
  \[
  \operatorname{Branch}(X_V/E_{\mathrm D})=D_b+D_\gamma,\qquad
  \operatorname{Branch}(Y/E_{\mathrm D})=D_a+D_\gamma,\qquad
  \operatorname{Branch}(X_H/E_{\mathrm D})=D_a+D_b.
  \]
\end{enumerate}
\end{conclusion}
\begin{proof}
由 \(\operatorname{Gal}(Z/E_{\mathrm D})\cong V_4\) 可知每个分支点惯性群均为阶 \(2\)。
记 \(N:=\deg\operatorname{Branch}(\pi)\)。因 \(g(E_{\mathrm D})=1,\ g(Z)=8\)，Hurwitz 公式给出
\[
2g(Z)-2
=\sum_{Q\in\operatorname{Branch}(\pi)}\frac{|V_4|}{|I_Q|}\bigl(|I_Q|-1\bigr)
=\sum_{Q\in\operatorname{Branch}(\pi)}2
=2N,
\]
故 \(14=2N\)，即 \(N=7\)。

设三条二次子覆盖对应非平凡特征
\(
\chi_V,\chi_Y,\chi_H\in V_4^\vee
\)。
对分支点 \(Q\)，记其惯性生成元为 \(s_Q\in V_4\setminus\{1\}\)。
二次子覆盖在 \(Q\) 处分歧当且仅当对应特征在 \(s_Q\) 上取值 \(-1\)。
对任意 \(s\neq 1\)，三条非平凡特征中恰有一条在 \(s\) 上取值 \(+1\)，其余两条取值 \(-1\)；
因此每个分支点处，三条二次子覆盖恰有一条不分歧。

记
\[
n_V:=\#\{Q\in\operatorname{Branch}(\pi):X_V\ \text{在}\ Q\ \text{处不分歧}\},
\]
\[
n_Y:=\#\{Q\in\operatorname{Branch}(\pi):Y\ \text{在}\ Q\ \text{处不分歧}\},\qquad
n_H:=\#\{Q\in\operatorname{Branch}(\pi):X_H\ \text{在}\ Q\ \text{处不分歧}\}.
\]
由结论 \ref{cor:xi-terminal-zm-s4-v4-three-quad-subcovers-branch-parity}，
\[
\deg\operatorname{Branch}(X_V/E_{\mathrm D})=6,\quad
\deg\operatorname{Branch}(Y/E_{\mathrm D})=6,\quad
\deg\operatorname{Branch}(X_H/E_{\mathrm D})=2.
\]
又 \(\#\operatorname{Branch}(\pi)=7\)，故
\[
7-n_V=6,\qquad 7-n_Y=6,\qquad 7-n_H=2,
\]
从而 \((n_V,n_Y,n_H)=(1,1,5)\)。
令 \(D_a,D_b,D_\gamma\) 分别为使 \(X_V\)、\(Y\)、\(X_H\) 不分歧的分支点集合，即得
\[
\deg D_a=1,\ \deg D_b=1,\ \deg D_\gamma=5
\]
及三条分支公式。证毕。
\end{proof}

\begin{conclusion}[第七分支点位于 \texorpdfstring{$E_{\mathrm D}\to\PP^1_y$}{ED to P1} 的分歧原像]\label{con:xi-terminal-zm-s4-v4-ed-seventh-branchpoint-cubic-ramified-preimage}
设 \(\varpi:E_{\mathrm D}\to\PP^1_y\) 为结论
\ref{thm:xi-s3-fiberproduct-branch-parity-xa-xv-ed} 中的三次覆盖，
其分歧值集合为
\[
\{0,1,\infty\}\cup\{g(y)=0\}.
\]
则：
\begin{enumerate}
  \item \(X_V\to E_{\mathrm D}\) 的 \(6\) 个分支点正是上述六个分歧值在 \(E_{\mathrm D}\) 上的未分歧原像点；
  \item 在结论 \ref{con:xi-terminal-zm-s4-v4-ed-branchcount-inertia-115} 中使 \(X_V\) 不分歧而 \(Z\) 分歧的唯一点，必为某个分歧值的分歧原像点。
\end{enumerate}
\end{conclusion}
\begin{proof}
第 \(1\) 点即结论 \ref{thm:xi-s3-fiberproduct-branch-parity-xa-xv-ed} 的第 \(3\) 点。

对第 \(2\) 点，记 \(S\subset\PP^1_y\) 为上述六个分歧值。
若 \(y_0\notin S\)，则 \(X\to\PP^1_y\) 在 \(y_0\) 上无分歧，故其上方点稳定子平凡。
任何中间商覆盖（特别是 \(Z\to E_{\mathrm D}\)）在该纤维上亦不可能新生分歧。
因此 \(\operatorname{Branch}(Z/E_{\mathrm D})\subset\varpi^{-1}(S)\)。

由结论 \ref{thm:xi-s3-fiberproduct-branch-parity-xa-xv-ed}，对每个 \(y_0\in S\)，纤维型为 \((2,1)\)：
记分歧原像为 \(R_{y_0}\)，未分歧原像为 \(U_{y_0}\)，且
\(
X_V\to E_{\mathrm D}
\)
恰在全部 \(U_{y_0}\) 处分歧。
由结论 \ref{con:xi-terminal-zm-s4-v4-ed-branchcount-inertia-115}，存在唯一分支点 \(P\) 使 \(X_V\) 在 \(P\) 处不分歧而 \(Z\) 分歧。
该点不能属于任何 \(U_{y_0}\)，故必有 \(P=R_{y_0}\) 对某个 \(y_0\in S\)。证毕。
\end{proof}

\begin{conclusion}[\texorpdfstring{$V_4$}{V4} 幂等关系导出的 Kani--Rosen 型分解]\label{con:xi-terminal-zm-s4-v4-ed-kanirosen-jacz-decomposition}
令 \(G:=\operatorname{Gal}(Z/E_{\mathrm D})\cong V_4\)，
并令 \(H_1,H_2,H_3\le G\) 为三条阶 \(2\) 子群，使
\[
Z/H_1\cong X_V,\qquad Z/H_2\cong Y,\qquad Z/H_3\cong X_H,\qquad Z/G\cong E_{\mathrm D}.
\]
则在 \(\QQ\)-同源意义下有
\[
\Jac(Z)\times \Jac(E_{\mathrm D})^{2}
\sim_{\QQ}
\Jac(X_V)\times \Jac(Y)\times \Jac(X_H).
\]
\end{conclusion}
\begin{proof}
在 \(\QQ[G]\) 中记
\[
e_H:=\frac{1}{|H|}\sum_{h\in H}h.
\]
对 \(G=V_4\) 的三条阶 \(2\) 子群直接计算得
\[
e_{H_1}+e_{H_2}+e_{H_3}=e_{\{1\}}+2e_G.
\]
将此恒等式作用于 \(\Jac(Z)\)。
按标准范数--拉回理论，
\(
\operatorname{Im}(e_H)\sim_{\QQ}\Jac(Z/H)
\) 且
\(
\operatorname{Im}(e_G)\sim_{\QQ}\Jac(Z/G)=\Jac(E_{\mathrm D})
\)。
于是得到
\[
\operatorname{Im}(e_{H_1})\times \operatorname{Im}(e_{H_2})\times \operatorname{Im}(e_{H_3})
\sim_{\QQ}
\Jac(Z)\times \Jac(E_{\mathrm D})^2,
\]
即所述分解。证毕。
\end{proof}

\begin{conclusion}[阿贝尔三维因子 \texorpdfstring{$B$}{B} 的局部 \texorpdfstring{$L$}{L} 因子有理消去公式]\label{con:xi-terminal-zm-s4-v4-ed-b-localfactor-rational-cancellation}
仍记
\(
B:=\Prym(Y/E_{\mathrm D})
\)。
若素数 \(p\) 处 \(Z,X_V,X_H,E_{\mathrm D}\) 皆良约化，记曲线 \(C\) 的局部 \(L\)-多项式分子为 \(P_p(C;T)\)，则
\[
P_p(B;T)
=
\frac{P_p(Z;T)\,P_p(E_{\mathrm D};T)}
{P_p(X_V;T)\,P_p(X_H;T)},
\qquad
\deg P_p(B;T)=6.
\]
等价地，在 Grothendieck 群中
\[
[H^1(B)]
=[H^1(Z)]+[H^1(E_{\mathrm D})]-[H^1(X_V)]-[H^1(X_H)],
\]
故对任意 \(n\ge 1\) 有
\[
\Tr\!\bigl(\mathrm{Frob}_p^n\mid H^1(B)\bigr)
=
\Tr\!\bigl(\mathrm{Frob}_p^n\mid H^1(Z)\bigr)
+\Tr\!\bigl(\mathrm{Frob}_p^n\mid H^1(E_{\mathrm D})\bigr)
-\Tr\!\bigl(\mathrm{Frob}_p^n\mid H^1(X_V)\bigr)
-\Tr\!\bigl(\mathrm{Frob}_p^n\mid H^1(X_H)\bigr).
\]
\end{conclusion}
\begin{proof}
由结论 \ref{con:xi-terminal-zm-s4-v4-ed-kanirosen-jacz-decomposition} 与
\ref{con:xi-terminal-zm-s4-c4-prym-threefold-and-v3sgn-identification}
中的
\(
\Jac(Y)\sim_{\QQ}E_{\mathrm D}\times B
\)，得到
\[
\Jac(Z)\times E_{\mathrm D}
\sim_{\QQ}
\Jac(X_V)\times \Jac(X_H)\times B.
\]
良素处同源阿贝尔簇具有相同 Frobenius 特征多项式，乘积对应局部因子相乘，故得比值公式。
再由 \(\dim B=3\)（同一结论已给出）得 \(\deg P_p(B;T)=2\dim B=6\)。
上式对 \(H^1\) 的 Grothendieck 群表达与迹恒等式与之等价。证毕。
\end{proof}

\begin{conclusion}[\texorpdfstring{$V_4$}{V4} 覆盖构造数据在椭圆基曲线上的刚化]\label{con:xi-terminal-zm-s4-v4-ed-abelian-cover-linebundle-rigidification}
设 \(D_a,D_b,D_\gamma\) 如结论
\ref{con:xi-terminal-zm-s4-v4-ed-branchcount-inertia-115}。
则存在线丛 \(L_a,L_b,L_\gamma\in\Pic(E_{\mathrm D})\)，满足
\[
L_a^{\otimes 2}\cong \mathcal O_{E_{\mathrm D}}(D_b+D_\gamma),\qquad
L_b^{\otimes 2}\cong \mathcal O_{E_{\mathrm D}}(D_a+D_\gamma),\qquad
L_\gamma^{\otimes 2}\cong \mathcal O_{E_{\mathrm D}}(D_a+D_b),
\]
且三条二次子覆盖分别由上述三组平方根数据给出。
特别地，
\[
\deg L_a=\deg L_b=3,\qquad \deg L_\gamma=1.
\]
此外 \(L_\gamma\in\Pic^1(E_{\mathrm D})\cong E_{\mathrm D}\) 等价于一点 \(Q\in E_{\mathrm D}\) 使
\[
2Q\sim D_a+D_b.
\]
\end{conclusion}
\begin{proof}
对光滑曲线上的 \(V_4\)-阿贝尔覆盖，标准构造数据为：分支除子
\(
\{D_s\}_{s\in V_4\setminus\{1\}}
\)
与角色线丛
\(
\{L_\chi\}_{\chi\in V_4^\vee\setminus\{1\}}
\)，满足
\[
L_\chi^{\otimes 2}\cong
\mathcal O_{E_{\mathrm D}}\!\Bigl(\sum_{\chi(s)=-1}D_s\Bigr).
\]
将 \(\chi_V,\chi_Y,\chi_H\) 与结论
\ref{con:xi-terminal-zm-s4-v4-ed-branchcount-inertia-115}
中的三条分支公式逐一对应，即得三条平方根关系。
次数由
\[
\deg L_\chi=\frac12\deg\!\Bigl(\sum_{\chi(s)=-1}D_s\Bigr)
\]
立即推出。
最后利用椭圆曲线上的标准同一
\(
\Pic^1(E_{\mathrm D})\cong E_{\mathrm D}
\)，即可将 \(L_\gamma\) 写成 \(2Q\sim D_a+D_b\)。证毕。
\end{proof}

\begin{conclusion}[三条 involution 商映射的分支极端不对称]\label{con:xi-terminal-zm-s4-v4-ed-involution-quotients-branch-2210}
令
\[
\iota_V,\iota_Y,\iota_H\in \operatorname{Gal}(Z/E_{\mathrm D})\setminus\{1\}
\]
分别对应
\[
Z/\langle\iota_V\rangle\cong X_V,\qquad
Z/\langle\iota_Y\rangle\cong Y,\qquad
Z/\langle\iota_H\rangle\cong X_H.
\]
则三条二次商映射分支次数为
\[
\deg\operatorname{Branch}(Z/X_V)=2,\qquad
\deg\operatorname{Branch}(Z/Y)=2,\qquad
\deg\operatorname{Branch}(Z/X_H)=10.
\]
等价地，
\[
\#\mathrm{Fix}_Z(\iota_V)=2,\qquad
\#\mathrm{Fix}_Z(\iota_Y)=2,\qquad
\#\mathrm{Fix}_Z(\iota_H)=10.
\]
\end{conclusion}
\begin{proof}
对任意二次覆盖 \(\rho:Z\to Q\) 有
\[
2g(Z)-2=2\bigl(2g(Q)-2\bigr)+\deg\operatorname{Branch}(\rho).
\]
代入 \(g(Z)=8\)、\(g(X_V)=4\)、\(g(Y)=4\)、\(g(X_H)=2\)，分别得
\[
14=2(6)+\deg\operatorname{Branch}(Z/X_V),\qquad
14=2(6)+\deg\operatorname{Branch}(Z/Y),
\]
\[
14=2(2)+\deg\operatorname{Branch}(Z/X_H),
\]
即分支次数依次为 \(2,2,10\)。
二次商中分支点与对应 involution 的固定点一一对应，故固定点个数结论等价成立。证毕。
\end{proof}

\begin{conclusion}[三条 Prym 因子的同源辨识与 \texorpdfstring{$B$}{B} 的唯一位置]\label{con:xi-terminal-zm-s4-v4-ed-three-prym-identifications}
定义
\[
P_V:=\Prym(Z/X_V),\qquad
P_Y:=\Prym(Z/Y),\qquad
P_H:=\Prym(Z/X_H).
\]
则在 \(\QQ\)-同源意义下
\[
P_V\sim_{\QQ}E\times B,
\]
\[
P_Y\sim_{\QQ}\Jac(X_A)\times E\times E_{\mathrm D},
\]
\[
P_H\sim_{\QQ}\Jac(X_A)\times E_{\mathrm D}\times B.
\]
\end{conclusion}
\begin{proof}
由结论 \ref{con:xi-terminal-zm-s4-delta-eigensplitting-jacz-prymxz}，
\[
\Jac(Z)\sim_{\QQ}\Jac(X_A)\times E\times E_{\mathrm D}^{2}\times B.
\]
另一方面：
\begin{itemize}
  \item 由定理 \ref{thm:xi-terminal-zm-s3-closure-jacobian-splitting}，
  \[
  \Jac(X_V)\sim_{\QQ}\Jac(X_A)\times E_{\mathrm D}^{2};
  \]
  \item 由结论 \ref{con:xi-terminal-zm-s4-c4-prym-threefold-and-v3sgn-identification}，
  \[
  \Jac(Y)\sim_{\QQ}E_{\mathrm D}\times B;
  \]
  \item 由结论 \ref{con:xi-terminal-zm-s4-xh-jacobian-splitting-ed-e}，
  \[
  \Jac(X_H)\sim_{\QQ}E\times E_{\mathrm D}.
  \]
\end{itemize}
对每条二次覆盖 \(Z\to Q\) 使用
\(
\Jac(Z)\sim_{\QQ}\Jac(Q)\times \Prym(Z/Q)
\)，并代入以上三式逐一消去公共因子，即得三条同源辨识。证毕。
\end{proof}

\endinput

\paragraph{\texorpdfstring{$(4,2^5)$}{(4,2^5)} 分歧型的 Nielsen 计数、辫单值群与 Jacobian 三挠几何}
以下结论在本节既定 \(S_4\)-正则覆盖记号下成立，并固定一个阶 \(4\) 分歧点为“无穷远位”。
记 \(\mathcal C_4\subset S_4\) 为四循环共轭类，\(\mathcal C_2\subset S_4\) 为换位共轭类。

\begin{conclusion}[内 Nielsen 类基数 \(160\) 与 \(B_5\) 传递性]\label{con:xi-terminal-zm-s4-nielsen-in-160-b5-transitive}
设
\[
\mathrm{Ni}^{\mathrm{in}}(S_4;\mathcal C_4,\mathcal C_2^5)
:=\left\{
(g_\infty,g_1,\ldots,g_5)\in S_4^6\ \middle|\
\begin{array}{l}
g_\infty\in \mathcal C_4,\ g_i\in\mathcal C_2,\\
g_\infty g_1\cdots g_5=1,\\
\langle g_\infty,g_1,\ldots,g_5\rangle=S_4
\end{array}
\right\}\Big/\!\sim,
\]
其中 \(\sim\) 为同时共轭。
则
\[
\bigl|\mathrm{Ni}^{\mathrm{in}}(S_4;\mathcal C_4,\mathcal C_2^5)\bigr|=160.
\]
并且 \(5\)-股辫群 \(B_5\)（仅作用于 \((g_1,\ldots,g_5)\)）在该 \(160\) 元集合上作用传递。
\end{conclusion}
\begin{proof}
设 Hurwitz 生成元 \(\sigma_i\ (1\le i\le 4)\) 作用为
\[
(g_i,g_{i+1})\longmapsto \bigl(g_{i+1},g_{i+1}^{-1}g_i g_{i+1}\bigr).
\]
脚本
\texttt{scripts/exp\_fold\_zm\_s4\_s3\_nielsen\_monodromy\_audit.py}
对 \(S_4\) 完全枚举给出原始满足约束的六元组数为 \(3840\)，按同时共轭商去后得到
\[
3840/|S_4|=3840/24=160.
\]
同一脚本对 \(\sigma_i,\sigma_i^{-1}\) 诱导图执行轨道遍历，证得唯一轨道，故 \(B_5\) 作用传递。证毕。
\end{proof}

\begin{conclusion}[\texorpdfstring{$S_4\twoheadrightarrow S_3$}{S4 to S3} 诱导的 \(40\) 块四重分块]\label{con:xi-terminal-zm-s4-to-s3-nielsen-40-blocks-of-4}
令 \(V_4^{\mathrm{norm}}\triangleleft S_4\) 为 Klein 四元正规子群，并取商映射
\[
\pi:S_4\longtwoheadrightarrow S_4/V_4^{\mathrm{norm}}\cong S_3.
\]
设 \(\mathcal T\subset S_3\) 为换位共轭类，则 \(\pi\) 诱导映射
\[
\pi_*:\mathrm{Ni}^{\mathrm{in}}(S_4;\mathcal C_4,\mathcal C_2^5)\longrightarrow
\mathrm{Ni}^{\mathrm{in}}(S_3;\mathcal T^6).
\]
并满足：
\begin{enumerate}
  \item \(\bigl|\mathrm{Ni}^{\mathrm{in}}(S_3;\mathcal T^6)\bigr|=40\)；
  \item \(\pi_*\) 满射，且每个纤维基数恒为 \(4\)；
  \item 上述 \(160\) 元集合因此分解为 \(40\) 个 \(B_5\)-不变块，每块大小 \(4\)。
\end{enumerate}
\end{conclusion}
\begin{proof}
同一脚本在 \(S_3\) 侧完全枚举得原始六元组总数 \(240\)，按同时共轭商去得到
\[
240/|S_3|=240/6=40.
\]
对每个 \(S_4\)-内类代表施行 \(\pi\) 并计数纤维，得到纤维大小分布恒为 \(4\)。
由 Hurwitz 作用与群同态、同时共轭均可交换，\(\pi_*\) 的纤维自动组成 \(B_5\)-不变分块系统。证毕。
\end{proof}

\begin{conclusion}[块上辫单值群识别为 \(\mathrm{PSp}_4(\mathbf F_3)\)]\label{con:xi-terminal-zm-braid-monodromy-psp4f3-identification}
记 \(\Gamma_{S_3}\subset S_{40}\) 为 \(B_5\) 在
\(\mathrm{Ni}^{\mathrm{in}}(S_3;\mathcal T^6)\) 上的像，\(\Gamma_{S_4}\subset S_{160}\) 为 \(B_5\) 在
\(\mathrm{Ni}^{\mathrm{in}}(S_4;\mathcal C_4,\mathcal C_2^5)\) 上的像。
则：
\begin{enumerate}
  \item \(|\Gamma_{S_3}|=25920\)，且 \(\Gamma_{S_3}\) 中心平凡、导群等于自身；
  \item 点稳定子次轨道分解为 \((1,12,27)\)；
  \item 因而 \(\Gamma_{S_3}\cong \mathrm{PSp}_4(\mathbf F_3)\)；
  \item \(160\) 点作用为带 \(40\) 个 \(4\)-块的非本原作用，其块上商像即 \(\Gamma_{S_3}\)。
\end{enumerate}
进一步地，若记
\[
|\Gamma_{S_4}|=2^{102}\cdot 3^{44}\cdot 5,
\]
则块核满足
\[
|K|=\frac{|\Gamma_{S_4}|}{|\Gamma_{S_3}|}=2^{96}\cdot 3^{40}.
\]
\end{conclusion}
\begin{proof}
脚本证书给出
\[
|\Gamma_{S_3}|=25920,\quad Z(\Gamma_{S_3})=1,\quad [\Gamma_{S_3},\Gamma_{S_3}]=\Gamma_{S_3},
\]
并给出次轨道 \((1,12,27)\)。
另一方面，\(\mathrm{PSp}_4(\mathbf F_3)\) 的阶为 \(25920\)，其在 \(\mathbf P^3(\mathbf F_3)\) 的 \(40\) 点作用为秩 \(3\)，次轨道正为 \((1,12,27)\)。
由这些群论不变量可识别块上单值群为 \(\mathrm{PSp}_4(\mathbf F_3)\)。

块结构由结论 \ref{con:xi-terminal-zm-s4-to-s3-nielsen-40-blocks-of-4} 给出，故 \(160\) 点作用非本原且块上商像为 \(\Gamma_{S_3}\)。
关于 \(|\Gamma_{S_4}|\) 与 \(|K|\) 的数值，脚本中保留了对应 Schreier--Sims 证书值及
\[
|\Gamma_{S_4}|/|\Gamma_{S_3}|=2^{96}3^{40}
\]
的一致性核验。证毕。
\end{proof}

\begin{conclusion}[四十个分块与属 \(2\) Jacobian 的循环 \(3\)-挠子群]\label{con:xi-terminal-zm-40-blocks-equal-cyclic-3torsion-subgroups}
设
\[
\varphi:X\to\PP^1
\]
为本节固定的 \(S_4\)-Galois 覆盖，分歧型 \((4,2,2,2,2,2)\)。
记
\[
C:=X/A_4,\qquad C_6:=X/V_4^{\mathrm{norm}}.
\]
则：
\begin{enumerate}
  \item \(C\to\PP^1\) 为在六个分歧值处分歧的二次覆盖，故 \(g(C)=2\)；
  \item \(C_6\to C\) 为次数 \(3\) 的循环无分歧覆盖；
  \item 该覆盖由 \(L\in \mathrm{Pic}^0(C)[3]\setminus\{0\}\)（等价于循环子群 \(\langle L\rangle\)）确定；
  \item 因 \(\mathrm{Pic}^0(C)[3]\cong (\ZZ/3\ZZ)^4\)，非平凡循环子群个数为
  \[
  \#\mathbf P^3(\mathbf F_3)=\frac{3^4-1}{3-1}=40;
  \]
  \item 结论 \ref{con:xi-terminal-zm-s4-to-s3-nielsen-40-blocks-of-4} 的 \(40\) 个块可几何识别为这些循环 \(3\)-挠子群。
\end{enumerate}
\end{conclusion}
\begin{proof}
第 \(1\) 点由 \(S_4/A_4\cong C_2\) 及结论
\ref{con:xi-terminal-zm-s4-normal-subgroup-quotients-explicit-models-genus}
直接得到。
第 \(2\) 点由 \(A_4/V_4^{\mathrm{norm}}\cong C_3\) 与结论
\ref{con:xi-terminal-zm-c6-over-c-etale-c3-pic3-class}
给出的无分歧性得到。
第 \(3\) 点为无分歧循环 \(3\)-覆盖与 \(3\)-挠线丛的标准对应。
第 \(4\) 点为有限线性代数计数。
第 \(5\) 点由结论 \ref{con:xi-terminal-zm-s4-to-s3-nielsen-40-blocks-of-4} 的基数与纤维分块，与上述 \(40\)-点参数空间同基数同构识别给出。证毕。
\end{proof}

\begin{theorem}[块轨道的辛正交刻画与强正则结构]\label{thm:xi-terminal-zm-block-orbits-symplectic-orthogonality-srg}
在结论 \ref{con:xi-terminal-zm-40-blocks-equal-cyclic-3torsion-subgroups} 的识别下，将
\[
\mathcal B:=\mathrm{Ni}^{\mathrm{in}}(S_3;\mathcal T^6)
\]
视为 \(V:=\mathrm{Pic}^0(C)[3]\cong (\ZZ/3\ZZ)^4\) 的一维 \(\FF_3\)-子空间（即循环子群）之集合。
记主极化所诱导的 Weil 配对为 \(e_3:V\times V\to\mu_3\)，并对两条线 \(\ell_1,\ell_2\subset V\) 定义关系
\[
\ell_1\perp \ell_2
\quad:\Longleftrightarrow\quad
e_3(x,y)=1\ \text{对任意}\ x\in\ell_1\setminus\{0\},\ y\in\ell_2\setminus\{0\}.
\]
则：
\begin{enumerate}
  \item 在 \(\mathcal B\) 上存在唯一的 \(\Gamma_{S_3}\)-不变非平凡对称二元关系，其在点稳定子下的两条非平凡子轨道基数为 \((12,27)\)；并且该关系与 \(\perp\) 完全一致。
  \item 因而以 \(\mathcal B\) 为顶点、以 \(\ell\neq\ell'\) 且 \(\ell\perp\ell'\) 为边的图为强正则图，参数为
  \[
  (v,k,\lambda,\mu)=(40,12,2,4).
  \]
  \item 其复置换表示分解满足
  \[
  \CC[\mathcal B]\cong \mathbf 1\oplus W_{24}\oplus W_{15},
  \qquad
  \dim W_{24}=24,\ \dim W_{15}=15,
  \]
  等价地邻接算子谱为 \(12,2,-4\)（重数分别为 \(1,24,15\)）。
\end{enumerate}
\end{theorem}

\begin{proof}
由结论 \ref{con:xi-terminal-zm-braid-monodromy-psp4f3-identification}，\(\Gamma_{S_3}\cong \mathrm{PSp}_4(\mathbf F_3)\) 在 \(\mathcal B\) 上为秩 \(3\) 的传递作用，点稳定子次轨道分解为 \((1,12,27)\)；因此任一 \(\Gamma_{S_3}\)-不变对称二元关系由两条非平凡子轨道之一唯一确定，从而至多存在一个非平凡关系。
\par
另一方面，在四维辛空间 \(V\cong\FF_3^4\) 的射影点集 \(\PP(V)\) 上，固定一条线 \(\ell\) 后其正交补 \(\ell^\perp\) 为三维子空间，故 \(\PP(\ell^\perp)\) 含
\(
(3^3-1)/(3-1)=13
\)
个点；去除 \(\ell\) 自身尚余 \(12\) 点，正给出 \(\ell\) 的正交邻域
\(
\mathcal B_\perp(\ell):=\{\ell'\neq\ell:\ \ell'\subset\ell^\perp\}
\)。
余下 \(27\) 点为非正交类。
因此 \(\perp\) 给出一条 \((12,27)\) 子轨道划分，并与上述唯一性一致，得到第 (1) 条。
\par
对强正则参数，若 \(\ell_1\perp\ell_2\) 且 \(\ell_1\neq\ell_2\)，则 \(\ell_1^\perp\cap \ell_2^\perp\) 为二维子空间，其射影点数为 \(4\)；其中包含 \(\ell_1,\ell_2\)，故公共邻点数为 \(2\)，即 \(\lambda=2\)。
若 \(\ell_1\not\perp\ell_2\)，则 \(\ell_1,\ell_2\not\subset \ell_1^\perp\cap\ell_2^\perp\)，公共邻点数为 \(4\)，即 \(\mu=4\)。
于是得到第 (2) 条。
第 (3) 条为强正则图的一般谱结论，与 Klingen Hecke 分解中的同一谱值一致（参见定理 \ref{thm:conclusion-m2-level3-delta0-inertia-klingen-charpoly}）。证毕。
\end{proof}

\begin{proposition}[由强正则图内生重构辛极空间 \(W(3)\) 的线系]\label{prop:xi-terminal-zm-reconstruct-W3-from-srg}
设图为定理 \ref{thm:xi-terminal-zm-block-orbits-symplectic-orthogonality-srg} 的强正则图。
则其所有极大团的大小均为 \(4\)。
记 \(\mathcal L\) 为全部 \(4\)-团之集合，则：
\begin{enumerate}
  \item \(|\mathcal L|=40\)；
  \item 每个顶点恰落在 \(4\) 个 \(4\)-团中；
  \item \((\mathcal B,\mathcal L,\in)\) 构成阶 \(3\) 的广义四边形（辛极空间）\(W(3)\)：每点在 \(4\) 条线上、每线含 \(4\) 点，且满足广义四边形公理。
\end{enumerate}
\end{proposition}

\begin{proof}
在 \(\mathcal B\simeq \PP(V)\) 的识别下，\(4\)-团等价于 \(\PP(M)\)，其中 \(M\subset V\) 为二维全各向同性子空间（Lagrangian 平面）。
确有 \(|\PP(M)|=(3^2-1)/(3-1)=4\)，且任意两点正交，故给出 \(4\)-团。
反向地，若四点两两正交，则其张成的子空间必为二维各向同性子空间，从而团必为某个 \(\PP(M)\)。
\par
计数：固定一点 \(\ell\in\PP(V)\)，包含 \(\ell\) 的 Lagrangian 平面个数等于 \(\PP(\ell^\perp/\ell)\) 的点数；而 \(\ell^\perp/\ell\) 为二维 \(\FF_3\)-空间，故该点数为 \((3^2-1)/(3-1)=4\)，得到第 (2) 条。
再由总旗数计数得 \(|\mathcal L|\cdot 4=|\mathcal B|\cdot 4=160\)，故 \(|\mathcal L|=40\)，得第 (1) 条。
广义四边形公理为辛极空间的标准性质：对点 \(\ell\) 与不含 \(\ell\) 的 Lagrangian 平面 \(M\)，存在唯一 \(\ell'\in\PP(M)\) 使 \(\ell\perp\ell'\)，从而唯一线 \(\langle \ell,\ell'\rangle\) 与 \(M\) 相交于 \(\ell'\)。证毕。
\end{proof}

\begin{theorem}[\texorpdfstring{\(\pi_*\)}{pi*} 的四重纤维与“通过一点的四条线”]\label{thm:xi-terminal-zm-pistar-fiber-equals-four-lines-through-point}
\leanverified{paper\_xi\_terminal\_zm\_pistar\_fiber\_equals\_four\_lines\_through\_point}
保持结论 \ref{con:xi-terminal-zm-s4-to-s3-nielsen-40-blocks-of-4} 的四重覆盖
\(
\pi_*:\mathrm{Ni}^{\mathrm{in}}(S_4;\mathcal C_4,\mathcal C_2^5)\to \mathcal B
\)。
固定 \(b\in\mathcal B\)，令 \(B:=\pi_*^{-1}(b)\) 为其纤维。
在定理 \ref{thm:xi-terminal-zm-block-orbits-symplectic-orthogonality-srg} 与命题 \ref{prop:xi-terminal-zm-reconstruct-W3-from-srg} 的识别下，将 \(b\) 视为 \(V\) 中一条线 \(\ell\subset V\)，并记
\[
\mathcal L(\ell):=\{M\in\mathcal L:\ \ell\in M\},
\qquad |\mathcal L(\ell)|=4.
\]
则存在一个由单值作用所确定的同构类，使得纤维 \(B\) 的点稳定子单值表示与 \(\mathcal L(\ell)\) 的自然作用同构；从而在该同构类意义下可将 \(B\) 识别为“通过点 \(\ell\) 的四条线”集合 \(\mathcal L(\ell)\)。
在此识别下，\(\ell\) 的正交邻域分裂为四个互不交叠的三元组
\[
\mathcal B_\perp(\ell)
=
\bigsqcup_{M\in\mathcal L(\ell)}\bigl(\PP(M)\setminus\{\ell\}\bigr),
\qquad |\mathcal B_\perp(\ell)|=12,
\]
并且上述分裂由纤维 \(B\) 内在标记。
\end{theorem}

\begin{proof}
由结论 \ref{con:xi-terminal-zm-braid-monodromy-psp4f3-identification}，块上单值群为 \(G=\Gamma_{S_3}\cong \mathrm{PSp}_4(\mathbf F_3)\)，并且 \(G\) 在 \(\mathcal B\) 上秩为 \(3\)。
设 \(P=\mathrm{Stab}_G(\ell)\)。
在辛几何中，集合 \(\mathcal L(\ell)\) 与 \(\PP(\ell^\perp/\ell)\cong \PP^1(\FF_3)\) 等价，且 \(P\) 在其上诱导的作用同构于 \(\mathrm{PGL}_2(\FF_3)\cong S_4\) 的自然四点作用，因此该四点传递作用在同构意义下唯一。
另一方面，\(\pi_*\) 给出大小为 \(4\) 的纤维，且纤维上的单值表示由点稳定子诱导；于是纤维单值表示与 \(\mathcal L(\ell)\) 的自然表示在同构类意义下必须一致。
最后的三元组分裂由 \(\PP(M)\) 各含四点且不同 \(M\) 仅交于 \(\ell\) 直接推出。证毕。
\end{proof}

\begin{corollary}[\texorpdfstring{\((3,3)\)}{(3,3)}-同源核的四重扩充律]\label{cor:xi-terminal-zm-cyclic3-to-33-isogeny-four-extensions}
令 \(A=\mathrm{Pic}^0(C)\) 为主极化阿贝尔曲面。
对任一非平凡循环子群 \(\ell\subset A[3]\)，包含 \(\ell\) 的极大各向同性子群 \(M\subset A[3]\) 恰有 \(4\) 个，并与 \(\mathcal L(\ell)\) 等势。
对每个 \(M\) 商 \(A/M\) 继承主极化，从而诱导一条 \((3,3)\)-同源
\[
\varphi_M:A\to A/M.
\]
因此，“给定循环 \(3\)-同源核 \(\ell\)”具有且仅具有四种方式可扩充为一个 \((3,3)\)-同源核 \(M\)。
\end{corollary}

\begin{proof}
极大各向同性二维子空间与 \((3,3)\)-同源核的对应由点--线四重纤维识别给出（见定理 \ref{thm:xi-terminal-zm-pistar-fiber-equals-four-lines-through-point}）。
而包含给定 \(\ell\) 的极大各向同性平面数目为 \(|\PP(\ell^\perp/\ell)|=4\)，与命题 \ref{prop:xi-terminal-zm-reconstruct-W3-from-srg} 一致。证毕。
\end{proof}

\begin{proposition}[线系的强正则对偶与旗空间的双 \(4\)-重投影]\label{prop:xi-terminal-zm-line-srg-and-flag-double-projection}
令 \(\mathcal L\) 为命题 \ref{prop:xi-terminal-zm-reconstruct-W3-from-srg} 的线系，并在 \(\mathcal L\) 上定义邻接关系
\[
M\sim N\quad:\Longleftrightarrow\quad \card{M\cap N}=3
\quad(\text{即 } M\cap N\ \text{为一条线}).
\]
则：
\begin{enumerate}
  \item \((\mathcal L,\sim)\) 亦为参数 \((40,12,2,4)\) 的强正则图；
  \item 点–线关联图 \(\mathcal I\subset \mathcal B\times\mathcal L\)（\((\ell,M)\in\mathcal I\iff \ell\subset M\)）为 \((4,4)\)-正则二部图，边数为 \(160\)；
  \item \(\mathcal I\to\mathcal B\) 与 \(\mathcal I\to\mathcal L\) 均为四重投影，并在组合意义下互为对偶：每个点邻域的 \(12\) 个正交邻点被唯一分裂为四个三元组，对应于通过该点的四条线；线邻域亦同理。
\end{enumerate}
\end{proposition}

\begin{proof}
在辛极空间 \(W(3)\) 中点与线完全对偶。
固定 \(M\in\mathcal L\)，其包含四条线 \(\ell\subset M\)；每条 \(\ell\) 被恰好四条线包含，其中之一为 \(M\) 本身，故与 \(M\) 交于 \(\ell\) 的其他线有三条。
因此 \(M\) 的邻接度为 \(4\cdot 3=12\)，余下为 \(27\)，与点图同参。
其余参数由同样的二维交与公共邻点计数得到，与点图的 \((\lambda,\mu)=(2,4)\) 完全平行。
二部图与双四重投影为命题 \ref{prop:xi-terminal-zm-reconstruct-W3-from-srg} 的计数结论之直接改写；对偶分裂由
\(
\mathcal B_\perp(\ell)=\bigsqcup_{M\in\mathcal L(\ell)}(\PP(M)\setminus\{\ell\})
\)
给出。证毕。
\end{proof}

\begin{conjecture}[内 Hurwitz 空间的抛物型 \(3\)-水平结构模解释]\label{conj:xi-terminal-zm-hurwitz-parabolic-level3-moduli}
设 \(\mathcal H_{S_3}^{\mathrm{in}}(\mathcal T^6)\) 与 \(\mathcal H_{S_4}^{\mathrm{in}}(\mathcal C_4,\mathcal C_2^5)\) 分别为相应分歧型别的内 Hurwitz 空间连通分量。
则期望存在一条与结论 \ref{con:xi-terminal-zm-40-blocks-equal-cyclic-3torsion-subgroups} 的 Jacobian 三挠识别与定理 \ref{thm:xi-terminal-zm-pistar-fiber-equals-four-lines-through-point} 的点–线–旗几何同时相容的双有理同构链，使得：
\begin{enumerate}
  \item \(\mathcal H_{S_3}^{\mathrm{in}}(\mathcal T^6)\) 与某个抛物型 \(3\)-水平结构的 Siegel 模叠（等价地：主极化阿贝尔曲面携带一条 \(3\)-挠线 \(\ell\subset A[3]\) 的模叠）在双有理意义下同构；
  \item 覆盖 \(\mathcal H_{S_4}^{\mathrm{in}}(\mathcal C_4,\mathcal C_2^5)\to \mathcal H_{S_3}^{\mathrm{in}}(\mathcal T^6)\) 对应于进一步选择一条包含 \(\ell\) 的极大各向同性子群 \(M\subset A[3]\)（即选择 \((3,3)\)-同源核），其几何纤维与 \(\PP(\ell^\perp/\ell)\cong \PP^1(\FF_3)\) 的四点截面一致。
\end{enumerate}
该猜想的可核验截面由 \(\Gamma_{S_3}\cong \mathrm{PSp}_4(\mathbf F_3)\) 的块上单值化、强正则结构与旗空间双投影给出；其几何提升需在代数簇层面构造相应的 \(A[3]\) 局部系统并验证其与 Hurwitz 单值数据的自然性一致。
\end{conjecture}

\begin{corollary}[Hurwitz 模空间不可约性与“属 \(2\)+\(3\)-挠数据”解释]\label{cor:xi-terminal-zm-hurwitz-irreducible-and-genus2-jac3torsion-moduli}
记 \(\mathcal H_{S_4}\) 为上述 \((4,2^5)\) 型 \(S_4\)-覆盖的 Hurwitz 模空间，
\(\mathcal H_{S_3}\) 为其 \(S_3\)-商型空间。则：
\begin{enumerate}
  \item \(\mathcal H_{S_4}\) 与 \(\mathcal H_{S_3}\) 均几何不可约，维数均为 \(3\)；
  \item 存在有限 \'{e}tale 映射 \(\mathcal H_{S_4}\to \mathcal H_{S_3}\)，泛纤维基数为 \(4\)；
  \item \(\mathcal H_{S_3}\) 可解释为“带六个分歧标记的属 \(2\) 曲线 \(C\) 及其 \(J(C)[3]\) 非平凡循环子群”之粗模空间。
\end{enumerate}
\end{corollary}
\begin{proof}
由结论 \ref{con:xi-terminal-zm-s4-nielsen-in-160-b5-transitive}、
\ref{con:xi-terminal-zm-s4-to-s3-nielsen-40-blocks-of-4} 的辫群传递性，Hurwitz 覆盖在分歧点配置空间上连通，从而几何不可约。
分歧点总数为 \(6\)，故维数为 \(6-3=3\)。
第 \(2\) 点由纤维大小恒为 \(4\) 得到。
第 \(3\) 点由结论 \ref{con:xi-terminal-zm-40-blocks-equal-cyclic-3torsion-subgroups} 给出的几何识别得到。证毕。
\end{proof}

\begin{proposition}[Hurwitz 生成元循环型的离散证书]\label{prop:xi-terminal-zm-hurwitz-generator-cycle-types-on-160-and-40}
在上述置换表示中，对任一标准生成元 \(\sigma_i\ (1\le i\le 4)\)：
\begin{enumerate}
  \item 在 \(160\) 点集合 \(\mathrm{Ni}^{\mathrm{in}}(S_4;\mathcal C_4,\mathcal C_2^5)\) 上，
  \[
  \sigma_i\sim 3^{36}2^{12}1^{28},
  \]
  因而置换阶为 \(6\)；
  \item 在 \(40\) 点集合 \(\mathrm{Ni}^{\mathrm{in}}(S_3;\mathcal T^6)\) 上，
  \[
  \sigma_i\sim 3^{9}1^{13},
  \]
  因而置换阶为 \(3\)。
\end{enumerate}
\end{proposition}
\begin{proof}
由脚本对四个 Hurwitz 生成元的循环分解直接给出：
\[
\{1^{28},2^{12},3^{36}\}\quad\text{与}\quad \{1^{13},3^{9}\},
\]
且四个生成元同型。置换阶由循环长度最小公倍数得到。证毕。
\end{proof}

\begin{conclusion}[模空间像的不可约性与一般点自同构群]\label{con:xi-terminal-zm-m16-locus-irreducible-generic-aut-s4}
令 \(\mathcal L\subset \mathcal M_{16}\) 为该 Hurwitz 家族在属 \(16\) 模空间中的像。
则：
\begin{enumerate}
  \item \(\mathcal L\) 为维数 \(3\) 的不可约代数簇；
  \item \(\mathcal L\) 一般点对应曲线 \(X\) 满足
  \[
  \mathrm{Aut}(X)=S_4.
  \]
\end{enumerate}
\end{conclusion}
\begin{proof}
第 \(1\) 点由推论
\ref{cor:xi-terminal-zm-hurwitz-irreducible-and-genus2-jac3torsion-moduli}
及模映射像的不可约性得到。
第 \(2\) 点中，对一般分歧配置，基底 \(\PP^1\) 上保持分歧集合的自同构群为平凡。
任何 \(X\) 的自同构都必须保持 \(S_4\)-覆盖结构并在基底诱导该平凡自同构，故仅剩甲板变换群 \(S_4\) 本身。证毕。
\end{proof}

\begin{remark}[可复现证书路径]\label{rem:xi-terminal-zm-s4-s3-nielsen-monodromy-certificate-path}
本组结论的离散计数与置换群证书由脚本
\texttt{scripts/exp\_fold\_zm\_s4\_s3\_nielsen\_monodromy\_audit.py}
生成，输出文件为
\texttt{artifacts/export/fold\_zm\_s4\_s3\_nielsen\_monodromy\_audit.json}。
\end{remark}

\endinput

